\section{After The Reader Switch}

The previous chapter ended by handing explanatory depth to the reader. That
handoff was not a retreat from the grammar. It was the last consequence of
the boundary reading. The interior could not choose its own dominant sign.
The cancelling pair required an exterior convention. The boundary restored a
true-not-false separator only after it could carry an on/off signal. The
reader then entered as the constrained role that orients that signal and
decides whether another explanation changes the record.

This chapter asks what such a reader is allowed to select. The answer is not
a private world, not an arbitrary convention, and not an absolute frame
standing above all frames. The reader selects a quotient. Presentations that
cannot be told apart by the admissible observation are identified. The
difference that survives that identification is read as relative motion,
orientation, charge, matter count, gauge partition, or predicate response,
depending on which grammatical door is being used. Relativity begins there:
not as a doctrine about looseness, but as the disciplined selection of what
remains after observational equivalence has been imposed.

The order matters. A frame is not given first and then used to manufacture
facts. A frame is the permitted selection by which a comparison becomes
readable. If the selection erases every difference between two
presentations, the grammar forbids the reader from spending that difference
as a new fact. If a residue survives, the grammar carries it. The residue is
relative because it belongs to the comparison; it is objective because the
comparison has specified what can be read and what must be quotiented away.

The native asymmetry of the title is the same discipline seen in sign and
count. Over the symmetric native baseline, the grammar reads matter and
counts no antimatter. That verdict is not a hidden preference for matter. It
is the count forced by the selected baseline. The positive sign appears only
when a tilted or asymmetric comparison is selected, and the grammar can then
read a closed-loop holonomy sign that the native open comparison did not
carry. Antimatter is not imported and then suppressed. It is unavailable
until the admissible comparison has supplied the orientation that can read
it.

The chapter therefore has five obligations. First, relative velocity must be
read as a quotient selection over observational equivalence. Second, the
positron must be read as a closed-loop holonomy sign, with open paths carrying
no charge. Third, the native matter/asymmetry count must be stated without
smuggling in an antimatter stock. Fourth, the finite gauge must partition
its thirty-six gates into pooling and selecting moves. Fifth, the reader must
become the last measuring station, the place where the predicates exposed by
the grammar turn back and ask whether a distinction has actually been read.

The continuity with the earlier chapters is strict. Difference is still
prior to sameness. Identity is still earned by quotient, anchor, or boundary.
Completion is still answerable to finite records. Signs are still
baseline-relative. The boundary still carries what the interior cannot spend.
What changes here is that the reader's selection is no longer a closing
afterthought. It becomes the grammar by which relativity, sign, asymmetry,
gauge, and final address are read together.
